\documentclass{ximera}
\title{Exponential functions}
\begin{document}
\begin{abstract}
The exponential function $b^x$, how it is computed for integer, rational and irrational exponents, and composing it with other functions. Textbook §1.4.
\end{abstract}
\maketitle
An exponential function is a function of the form
\[
f(x)= \answer{b^{x}}
\]
where $b$ is a positive constant (i.e., $b \in \mathbb{R}^{+}$).

How do we calculate $b^{x}$ for every $x \in \mathbb{R}$?

% AUTHOR CHOICE: the "Example" column was left empty in the original; examples chosen during conversion, change freely
% work: 2^3 = 8; 8^{2/3} = (cube root of 8)^2 = 2^2 = 4
\begin{center}
\begin{tabular}[t]{|l|l|l|}
\hline
$x$ & $b^{x}$ & Example \\
\hline
Integer ( $\mathbb{Z}$ ) & $b^{x}$ & $2^{3}=\answer{8}$ \\
\hline
Rational $(\mathbb{Q})$ & $b^{p / q}=\sqrt[q]{b^{p}}$ & $8^{2/3}=\answer{4}$ \\
\hline
Irrational $(\mathbb{R} \backslash \mathbb{Q})$ & Estimate using graph & $2^{\sqrt{2}}$ \\
\hline
\end{tabular}
\end{center}

\begin{example}\label{example:1-10}
% AUTHOR CHOICE: example chosen during conversion, change freely
Let $f(x)=2^{x}$ and let $g(x)=x+1$.
\[
(f \circ g)(x) = \answer{2^{x+1}} \qquad (g \circ f)(x) = \answer{2^{x}+1}
\]
Are these the same function? (i.e., is $f \circ g=g \circ f$) \wordChoice{\choice{Yes}\choice[correct]{No}}.
\end{example}
\end{document}
